现在注意到态射
\begin{equation}\tag{3.6.13}
\left\{
\begin{aligned}
\psi &:V'\longrightarrow
  \mathbb P^r\times V'\subset\mathbb P^r\times\check X,\\
\psi(v)&=(\psi_1(v),v)
\end{aligned}
\right.
\end{equation}
可延拓为态射
\begin{equation}\tag{3.6.14}
\left\{
\begin{aligned}
\widetilde\psi &:V\longrightarrow\mathbb P^r\times V,\\
\widetilde\psi(v)&=(\widetilde\psi_1(v),v),
\end{aligned}
\right.
\end{equation}
其中
\begin{equation}\tag{3.6.15}
\widetilde\psi_1:V\longrightarrow\mathbb P^r
\end{equation}
为 Gauss 态射, 现在将它视作定义在整个 $V$ 上.

我们刚证明了
\begin{equation}\tag{3.6.16}
\psi(V')\subset\mathbb P(N).
\end{equation}
由于 $V'$ 在 $V$ 中稠密, 且 $\mathbb P(N)$ 在 $\mathbb P^r\times\check X$ 中闭, 由此可得
\begin{equation}\tag{3.6.17}
\widetilde\psi(V)\subset\mathbb P(N),
\end{equation}
而且显然有
\begin{equation}\tag{3.6.18}
\widetilde\psi(V)\subset\varphi^{-1}(V).
\end{equation}
再由连续性, 恒等式
\[
\psi\circ\varphi=\operatorname{id}_{\varphi^{-1}(V')},
\qquad
\varphi\circ\psi=\operatorname{id}_{V'}
\]
延拓为
\[
\widetilde\psi\circ\varphi=\operatorname{id}_{\varphi^{-1}(V)},
\qquad
\varphi\circ\widetilde\psi=\operatorname{id}_{V}.
\]
所以 3.5(a) 和 3.5(b) 中的开子集相同.

接着证明 3.5(a) 和 3.5(c) 中的开子集相同. 由 3.3, 后者就是由
\begin{equation}\tag{3.6.19}
U=\left\{H\in\check X\ \middle|\
\begin{array}{l}
\varphi^{-1}(H)\text{ 恰由一个点组成, 且}\\
\varphi\text{ 在该点处非分歧}
\end{array}
\right\}.
\end{equation}
定义的子集 $U\subset\check X$. 概形 $\varphi^{-1}(U)$ 光滑, 因而正规; 诱导态射
\[
\varphi:\varphi^{-1}(U)\longrightarrow U
\]
由已证结果为双有理态射, 由 $U$ 的定义为根式态射, 并且固有. 所以它是 $U$ 的正规化. 特别地, $U$ 几何单支. 由于此态射非分歧, EGA IV, 18.10.1 表明它平展; 再由 EGA IV, 18.10.18 可知它是开浸入. 因而它是同构, 且 $U$ 光滑. 这证明了 $3.5(c)\subset 3.5(a)$. 公式 (3.6.19) 使包含关系 $3.5(b)\subset 3.5(c)$ 显然, 再由等式 $3.5(a)=3.5(b)$ 完成证明.